\documentclass[12pt,a4paper]{article}
\usepackage[margin=25mm, headheight=15pt, headsep=10pt]{geometry}
\usepackage{fontspec}
\usepackage[table]{xcolor} % Note: table option is required for rowcolors
\usepackage{titlesec}
\usepackage{tocloft}
\usepackage{fancyhdr}
\usepackage{booktabs}
\usepackage{tcolorbox}
\tcbuselibrary{skins, breakable}
\usepackage[colorlinks=true, linkcolor=emerald, urlcolor=emerald, bookmarks=true]{hyperref}

\definecolor{emerald}{RGB}{0,102,51}
\definecolor{darkred}{RGB}{139,0,0}
\definecolor{lightgray}{RGB}{245,245,245}

\usepackage[nil]{babel}
\babelprovide[import, main]{bengali}
\babelprovide[import]{arabic}
\babelfont{rm}[Renderer=HarfBuzz,Scale=1.0]{Noto Serif Bengali}
\babelfont{sf}[Renderer=HarfBuzz]{Noto Sans Bengali}
\babelfont{tt}[Renderer=HarfBuzz]{Noto Sans Bengali}
\babelfont[arabic]{rm}[Renderer=HarfBuzz,Scale=1.1]{Noto Naskh Arabic}

% Section styling
\titleformat{\section}{\Large\bfseries\color{emerald}}{\thesection.}{0.5em}{}
\titleformat{\subsection}{\large\bfseries\color{emerald!85!black}}{\thesubsection.}{0.5em}{}
\titleformat{\subsubsection}{\normalsize\bfseries\color{emerald!70!black}}{\thesubsubsection.}{0.5em}{}
\titlespacing*{\section}{0pt}{18pt}{8pt}

% Fancy Header/Footer setup
\pagestyle{fancy}
\fancyhf{} % clear all header and footer fields
\fancyhead[L]{\color{emerald}\small\textbf{\nouppercase{\leftmark}}}
\fancyfoot[L]{\scriptsize\color{black!50}{\lat{AI}}-সংকলিত, আলেম-যাচাইকৃত নয়}
\fancyfoot[C]{\color{emerald!70!black}\thepage}
\renewcommand{\headrulewidth}{0.4pt}
\renewcommand{\headrule}{\hbox to\headwidth{\color{emerald}\leaders\hrule height \headrulewidth\hfill}}
\renewcommand{\footrulewidth}{0pt}

% Custom boxes
% 1. Ayat/Hadith Box
\newtcolorbox{ayatbox}[1][]{
    enhanced, breakable, 
    colback=emerald!5, 
    colframe=emerald, 
    boxrule=0.5pt, 
    arc=3pt, 
    left=10pt, right=10pt, top=10pt, bottom=10pt,
    #1
}

% 2. Quote/Translation Box
\newtcolorbox{quotebox}[1][]{
    enhanced, breakable,
    frame hidden,
    colback=lightgray,
    borderline west={3pt}{0pt}{emerald},
    left=15pt, right=10pt, top=10pt, bottom=10pt,
    sharp corners,
    #1
}

% 3. AI-generated notice box
\newtcolorbox{warnbox}[1][]{
    enhanced, breakable,
    colback=red!4,
    colframe=darkred,
    boxrule=0.8pt,
    arc=3pt,
    left=10pt, right=10pt, top=10pt, bottom=10pt,
    #1
}

% Commands
\newcommand{\arab}[1]{\foreignlanguage{arabic}{#1}}
\newcommand{\kib}[1]{\textcolor{emerald}{\textbf{#1}}}

% Latin (English) fallback: Noto Serif Bengali has NO Latin glyphs
\newfontfamily\latfont[Renderer=HarfBuzz]{Noto Serif}
\newcommand{\lat}[1]{{\latfont #1}}

\setlength{\parskip}{5pt}
\setlength{\parindent}{0pt}

% Fix for missing bullet in Noto Serif Bengali
\renewcommand{\labelitemi}{$\bullet$}



\begin{document}
\begin{center}{\Large\bfseries\color{emerald} আল্লাহর রহমত ও তাওবা — \lat{11}-যাদহানের-তাওবা-ইবনে-মাসউদের-হাতে}\end{center}
\vspace{1em}
\section*{ঘটনা ১১ — তাবেয়ী যাদহান: এক কঠোর স্মরণ করিয়ে দেওয়ায় গোটা জীবন বদলে যাওয়া (সালাফের ঘটনা)}

\begin{warnbox}
 \textbf{গুরুত্বপূর্ণ নোটিশ (\lat{Important} \lat{Notice}):} এই কন্টেন্টটি \lat{AI} (কৃত্রিম বুদ্ধিমত্তা) দিয়ে প্রস্তুত ও সংকলিত। \lat{AI} কঠোর সূত্র-মান নীতি মেনে শুধু নির্ভরযোগ্য উৎস থেকে তথ্য সংগ্রহ করেছে — কেবল সহীহ/হাসান হাদিস ও সহীহ সনদের আসার রাখা হয়েছে, দুর্বল সনদ বাদ দেওয়া হয়েছে। তবে এটি কোনো আলেম বা মুহাদ্দিস কর্তৃক যাচাইকৃত নয়।
\end{warnbox}

\begin{quote}
\textbf{মান-সংক্রান্ত সৎ ঘোষণা:} এটি নবীজির (সা.) মারফু' হাদিস নয় — এটি \textbf{তাবেয়ী যাদহানের নিজের ঘটনার বর্ণনা} (আসার), যা ইমাম যাহাবী তাঁর সিয়ারু আলাম আন-নুবালায় যাদহানের জীবনীতে এনেছেন। সনদের একজন বর্ণনাকারী (আবু হাশিম আল-রুম্মানি) নিয়ে রিজালবিদদের মতভেদ আছে (কেউ কেউ দুর্বল বলেছেন) — তাই একে \textbf{সহীহ হাদিসের মর্যাদা দেওয়া যায় না}; সালাফের নির্ভরযোগ্য জীবনী-সাহিত্যের বর্ণনা হিসেবে উপস্থাপন করা হলো।
\end{quote}

\medskip\hrule\medskip

\subsection*{১. সূত্র কার্ড (\lat{Source} \lat{Card})}

\begin{table}[h]\centering\small
\begin{tabular}{ll}
বিষয় & বিবরণ \\
\hline
\textbf{ঘটনা} & যুবক যাদহান বাদ্যযন্ত্র ও নবীয (খেজুরের নেশা-কর পানীয়) নিয়ে গান গাইছিল; আবদুল্লাহ ইবনে মাসউদ (রাঃ) এসে সেগুলো ভেঙে দিয়ে বললেন — তোমার সুন্দর কণ্ঠ যদি কুরআন তিলাওয়াতে লাগত! এই কথায় যাদহানের অন্তরে তাওবা জাগল — সে কাঁদতে কাঁদতে ইবনে মাসউদের পিছু নিল, আর তার জীবন বদলে গেল \\
\textbf{বর্ণনাকারী} & আবু হাশিম আল-রুম্মানি $\leftarrow$ যাদহান (নিজের ঘটনা) \\
\textbf{গ্রন্থ} & সিয়ারু আলাম আন-নুবালা (ইমাম যাহাবী) — যাদহানের জীবনী \\
\textbf{সনদের মান} & \textbf{জীবনীগ্রন্থের বর্ণনা (আসার) — মারফু' হাদিস নয়;} সনদে মতভেদ আছে (বর্ণনাকারী আবু হাশিম আল-রুম্মানিকে কেউ কেউ দুর্বল বলেছেন) \\
\textbf{স্থান ও সময়} & কুফা — তাবেয়ী যুগ; যাদহান হিজরি ৮২ সালে ইন্তেকাল করেন \\
\hline
\end{tabular}
\end{table}

\medskip\hrule\medskip

\subsection*{২. সংক্ষিপ্ত পটভূমি (\lat{Background})}

যাদহান (রহ.) ছিলেন তাবেয়ী যুগের (ইসলামের দ্বিতীয় প্রজন্ম) একজন মানুষ — নবীজির (সা.) যুগে জন্মেছিলেন, কিন্তু নবীজিকে পাননি। তার বয়সের অন্য অনেকের মতো তিনিও \textbf{ইবনে মাসউদের (রাঃ) ছাত্রদের একজন} হয়ে ওঠেন। কিন্তু তার শুরুটা ছিল অন্যরকম — যৌবনে তিনি বাদ্যযন্ত্র (\lat{tanbur}) বাজাতেন, নবীয পান করতেন এবং বন্ধুদের গান শোনাতেন। একদিন আবদুল্লাহ ইবনে মাসউদ (রাঃ) — যিনি কুফায় ছিলেন এবং যাঁর কাছ থেকে যাদহান হাদিসও বর্ণনা করেছেন — তার কাছে দিয়ে যাচ্ছিলেন। যা ঘটল, তা যাদহানের জীবনের মোড় ঘুরিয়ে দিল।

\medskip\hrule\medskip

\subsection*{৩. পূর্ণ ঘটনার বিশুদ্ধ বর্ণনা (\lat{The} \lat{Full} \lat{Narration})}

যাদহান (রহ.) নিজেই বর্ণনা করেন:

\begin{quote}
\lat{“}আমি ছিলাম সুন্দর কণ্ঠের এক যুবক — বাদ্যযন্ত্র (\lat{tanbur}) বাজাতে পারদর্শী। একবার আমি এক বন্ধুর কাছে ছিলাম — আমাদের কাছে নবীয (খেজুরের পানীয়) ছিল, আর আমি তাদের গান শুনাচ্ছিলাম। এমন সময় আবদুল্লাহ ইবনে মাসউদ (রাঃ) পাশ দিয়ে যাচ্ছিলেন — তিনি আমাদের কাছে ঢুকে পড়লেন, \textbf{নবীযের পাত্রটিতে আঘাত করলেন এবং বাদ্যযন্ত্রটি ভেঙে দিলেন।} তারপর বললেন: \textbf{\lat{“}ওহে বালক! তোমার এই সুন্দর কণ্ঠ যদি কুরআন তিলাওয়াতে ব্যবহৃত হতো — তবে তুমিই হতে সেই ব্যক্তি!\lat{”}} এই বলে তিনি চলে গেলেন।
\end{quote}

\begin{quote}
আমি আমার বন্ধুদের জিজ্ঞেস করলাম: \lat{“}ইনি কে?\lat{”} তারা বলল: \textbf{\lat{“}ইনি আবদুল্লাহ ইবনে মাসউদ।\lat{”}} — তখনই আমার অন্তরে তাওবার আলো (\lat{tawbah}) জাগল। আমি \textbf{কাঁদতে কাঁদতে তার পিছু ধাবিত হলাম}, তার কাপড় ধরে ফেললাম। তিনি আমার দিকে ফিরলেন — আমাকে জড়িয়ে ধরে তিনিও কাঁদলেন। তারপর বললেন: \textbf{\lat{“}স্বাগতম তার জন্য, যাকে আল্লাহ ভালোবাসেন। বসো।\lat{”}} তিনি ভেতরে গিয়ে আমার জন্য কিছু খেজুর নিয়ে এলেন।\lat{”}
\end{quote}

এরপর যাদহান (রহ.) হয়ে ওঠেন নিষ্ঠাবান ইবাদতকারী। ইমাম যাহাবী (রহ.) বলেন — তাকে দেখা যেত নামাজে এমন নিশ্চল হয়ে দাঁড়িয়ে থাকতে, যেন তিনি একটি কাঠের টুকরো বা গাছ। তিনি হিজরি ৮২ সালে ইন্তেকাল করেন।

(সিয়ারু আলাম আন-নুবালা — যাদহানের জীবনী)

\medskip\hrule\medskip

\subsection*{৪. শব্দের অর্থ (\lat{Key} \lat{Words})}

\begin{table}[h]\centering\small
\begin{tabular}{ll}
শব্দ & অর্থ \\
\hline
\textbf{আত-তানবুর (\arab{الطنبور})} & তার-যুক্ত বাদ্যযন্ত্র — গানের অনুষঙ্গ \\
\textbf{আন-নবীয (\arab{النبيذ})} & খেজুর-ভেজানো পানীয় — সময়ের ব্যবধানে নেশাকর হতে পারে \\
\textbf{তাওবা (\arab{توبة})} & ফিরে আসা — পাপ থেকে আল্লাহর দিকে \\
\textbf{আল-খুশু' (\arab{الخشوع})} & নামাজে নিশ্চল বিনয় — অন্তরের উপস্থিতি \\
\hline
\end{tabular}
\end{table}

\medskip\hrule\medskip

\subsection*{৫. সংশ্লিষ্ট দলিল ও রেফারেন্স}

\begin{itemize}
\item \textbf{আয়াত (আল-ফুরকান ২৫:৭০):} তাওবার পর পাপ নেকিতে বদল — যাদহানের (রহ.) জীবনে এর বাস্তব রূপ দেখা যায়: গান-বাজনার জায়গায় এসেছে কুরআন তিলাওয়াত ও নামাজের খুশু'।
\item \textbf{হাদিস (মুসলিম ২৭৫৯):} \lat{“}আল্লাহ রাতের বেলায় হাত বাড়ান, যেন দিনের পাপীরা তাওবা করে; আর দিনে হাত বাড়ান, যেন রাতের পাপীরা তাওবা করে…\lat{”} — তাওবার দরজা সবসময় খোলা।
\item \textbf{হাদিস (বুখারি ৬৩০৮; মুসলিম ২৭৪৪):} তাওবায় আল্লাহর খুশি — হারানো উট ফেরত পাওয়ার সাদৃশ্য।
\item \textbf{ইবনে মাসউদের (রাঃ) পদ্ধতি:} মন্দ কাজে লিপ্ত যুবককে তাড়ানো নয় — বরং তার ভালো গুণ (কণ্ঠ) ধরিয়ে দেওয়া, যেন সে তা আল্লাহর কাজে লাগায়। এটি নববী দাওয়াহ পদ্ধতির তাবেয়ী যুগের নমুনা।
\end{itemize}
\subsubsection*{মূল শিক্ষা (দলিলভিত্তিক)}

\begin{enumerate}
\item \textbf{একটি সত্য কথা জীবন বদলে দিতে পারে:} ইবনে মাসউদের (রাঃ) কয়েকটি কথা যাদহানের (রহ.) অন্তরে তাওবা জাগাল। নসীহতের (\lat{sincere} \lat{advice}) শক্তি অনুমান করা যায় না — আল্লাহ তাওবাকারীর অন্তরে সেই আলো নিজেই জাগিয়ে দেন।
\item \textbf{পরিবেশ বদলানো জরুরি:} যাদহান (রহ.) গান-বাজনা ও নবীযের আসর ছেড়ে ইলম ও ইবাদতের আসরে গেলেন — তাওবার পর পরিবেশ (\lat{environment}) বদলানো তাওবার স্থায়িত্বের জন্য অপরিহার্য।
\item \textbf{আলেমের স্নেহ ও দোয়া:} ইবনে মাসউদ (রাঃ) যুবকটিকে ধমক দিয়ে তাড়াননি — জড়িয়ে ধরে কেঁদেছেন, খেজুর এনে দিয়েছেন। পাপীকে ফেরানোর জন্য কঠোরতা নয় — ভালোবাসা ও কাছে টানা।
\item \textbf{একজন সালাফের ঘটনার শিক্ষা:} যাদহান (রহ.) নবীজিকে (সা.) পাননি, তবু তাঁর কণ্ঠ ও জীবন আল্লাহর কাজে লেগেছে — ইসলামের ইতিহাসে (\lat{history}) এমন অসংখ্য মানুষ আছেন যারা যেকোনো বয়সে ফিরে এসে বড় আলেম-আবিদ হয়েছেন।
\end{enumerate}
\medskip\hrule\medskip

\subsection*{৬. রেফারেন্স}

\begin{itemize}
\item সিয়ারু আলাম আন-নুবালা (ইমাম যাহাবী) — যাদহানের জীবনী
\item সূত্র: আবু হাশিম আল-রুম্মানি $\leftarrow$ যাদহান
\item সনদ-টীকা: আবু হাশিম আল-রুম্মানির নির্ভরযোগ্যতা নিয়ে মতভেদ আছে — আল-খাযরাজির খুলাসাতুত তাহযিবে তাকে \lat{“}মাজহুল/দুর্বল\lat{”} বলা হয়েছে; তাই এই বর্ণনা সহীহ হাদিসের মর্যাদা পায় না
\item কুরআন: আল-ফুরকান ২৫:৭০; আয-যুমার ৩৯:৫৩
\item হাদিস: মুসলিম ২৭৫৯, ২৭৪৪; বুখারি ৬৩০৮
\end{itemize}

\end{document}
